\section{Continuous nowhere-differentiable functions}
\label{sec:nowhere-differentiable}
No point of a nowhere differentiable function has a finite derivative.
Such a function can nevertheless be continuous. Section~38
constructs examples with both properties.
When an exercise needs such a function, Theorem~38.1 supplies
a bounded continuous example on $\R$.

One useful strategy is to suppose a derivative exists and use
derivative rules to force a derivative of the known nowhere-differentiable
function. This is a contradiction argument: the rules are applied
only to functions that are differentiable under the temporary
assumption, never directly to a function already known to be
nondifferentiable.

In the metric-space notation used in this section, $C(\R)$ consists
of \emph{bounded} continuous functions, and
$d(f,g)=\sup_{x\in\R}|f(x)-g(x)|$. The notation $C^\infty$
instead means that derivatives of every finite order exist and are
continuous. These are two different uses of the letter $C$.

\exercise{38.1}{Why a Taylor series cannot produce the desired example}
\goal Explain the obstruction to constructing a continuous
nowhere-differentiable function as a convergent Taylor series.
\step{Recall what a positive-radius power series guarantees.}
Suppose
\[
 f(x)=\sum_{n=0}^\infty a_n(x-c)^n
\]
has radius of convergence $R>0$. On the open interval
$(c-R,c+R)$, the power-series differentiation theorem gives
\[
 f'(x)=\sum_{n=1}^\infty n a_n(x-c)^{n-1}.
\]
In particular $f$ is differentiable at every interior point of its
convergence interval, however complicated the coefficients might be.
Repeating the theorem gives derivatives of every order there.
\step{Check the possible loopholes.}
Behavior at the two endpoints of a finite convergence interval
can be different, but those endpoints do not remove differentiability
at all the interior points. A radius of convergence $R=0$ gives
only a value at the center, not a function on an interval.
Moreover, a literal Taylor series of a function requires derivatives
at its center in order to define the coefficients in the first place.
\answer{A positive-radius Taylor series is necessarily differentiable
on its interior, so it cannot be a nowhere-differentiable function.}

\exercise{38.2}{A polynomial derivative forces a polynomial function}
\goal If $f^{(n)}$ is a polynomial on $\R$ for some nonnegative
integer $n$, prove that $f$ is a polynomial.
\step{Build a polynomial with the prescribed $n$th derivative.}
The case $n=0$ already says that $f$ is a polynomial. For $n\geq1$,
write
\[
 f^{(n)}(x)=p(x)=\sum_{k=0}^m a_kx^k,
 \qquad
 P(x)=\sum_{k=0}^m a_k\frac{k!}{(k+n)!}x^{k+n}.
\]
Differentiating $x^{k+n}$ exactly $n$ times multiplies it by
$(k+n)!/k!$ and leaves $x^k$. Hence $P^{(n)}=p$.
\step{Subtract and use the zero-derivative principle.}
Let $h=f-P$. Then $h^{(n)}=0$ everywhere. Proposition~38.4
implies that $h$ is a polynomial of degree at most $n-1$.
To recall why: $h^{(n-1)}$ has derivative zero, so the Mean Value
Theorem makes it constant. Subtract an antiderivative of that
constant from $h^{(n-2)}$; the resulting derivative is zero,
so that difference is constant as well. Repeating gives a polynomial
expression for $h$ after $n$ steps.
\answer{$f=P+h$ is a sum of polynomials and is therefore a polynomial.}

\exercise{38.3}{A differentiable function with a nowhere-differentiable derivative}
\goal Construct a differentiable function whose first derivative
is nowhere differentiable.
\step{Integrate a continuous nowhere-differentiable function.}
Let $w$ be a continuous nowhere-differentiable function on $\R$,
whose existence is proved in this section. Define
\[
 H(x)=\int_0^x w(t)\,dt.
\]
For every real $x$, this is an ordinary integral over a finite
interval (with reversed orientation if $x<0$), so it exists.
\step{Use the hypothesis that makes the Fundamental Theorem applicable.}
Since $w$ is continuous at every $x$, the Fundamental Theorem of
Calculus gives $H'(x)=w(x)$ at every real $x$. Thus $H$ is
differentiable everywhere, and $H'$ is nowhere differentiable.
Equivalently, $H$ has no second derivative at any point.
\answer{An antiderivative of a continuous nowhere-differentiable
function is the required example.}

\exercise{38.4}{Adding a differentiable function cannot repair a derivative}
\goal If $f$ is nowhere differentiable and $g$ is differentiable
on $\R$, decide whether $f+g$ must be nowhere differentiable.
\step{Suppose there were an exceptional point.}
Assume, for contradiction, that $f+g$ is differentiable at some
$a$. By hypothesis $g$ is differentiable at that same $a$.
The difference of two functions differentiable at a point is
differentiable there. Thus
\[
 f=(f+g)-g
\]
would be differentiable at $a$, with derivative
$(f+g)'(a)-g'(a)$. This contradicts nowhere differentiability of $f$.
\answer{Yes. The sum $f+g$ is nowhere differentiable.}

\exercise{38.5}{Two rough functions can cancel in a sum or a product}
\step{(a) A sum need not remain nowhere differentiable.}
Take any continuous nowhere-differentiable function $w$ and let
$f=w$, $g=-w$. The function $g$ is also nowhere differentiable:
if $g'(a)$ existed, multiplying by $-1$ would give $w'(a)$.
Both are continuous, but $f+g=0$, which is differentiable everywhere.
\step{(b) Arrange a nonzero function before taking its reciprocal.}
Choose a bounded continuous nowhere-differentiable $w$ and a finite
bound $M\geq0$ with $|w(x)|\leq M$. Define
\[
 f(x)=1+\frac{w(x)}{2(M+1)}.
\]
This is continuous, satisfies
\[
 d(f,1^*)\leq\frac{M}{2(M+1)}<\frac12<1,
 \qquad \frac12<f(x)<\frac32,
\]
and is nowhere differentiable: differentiability of $f$ would imply
differentiability of $w=2(M+1)(f-1)$ at the same point.
\step{Verify that the reciprocal is itself nowhere differentiable.}
Let $g(x)=1/f(x)$. It is continuous because $f$ never vanishes.
Suppose $g$ were differentiable at a point $a$. We have
$g(a)=1/f(a)\ne0$. The reciprocal differentiation rule, applied
to this temporarily differentiable function $g$, would then make
$1/g=f$ differentiable at $a$, with derivative
$-g'(a)/g(a)^2$. This contradicts the established property of $f$.
Thus $g$ is nowhere differentiable too, yet $f(x)g(x)=1$ everywhere.
\answer{Neither a sum nor a product of two continuous
nowhere-differentiable functions must be nowhere differentiable.}

\exercise{38.6}{A quotient of smooth functions is smooth}
\goal Suppose $f,g\in C^\infty(\R)$ and $g(x)\ne0$ for every $x$.
Prove $f/g\in C^\infty(\R)$.
\idea Inductively describe each derivative so that its denominator
is a power of the nonvanishing function $g$.
\step{State a formula suited to induction.}
We claim that, for every integer $n\geq0$, there is
$h_n\in C^\infty(\R)$ for which
\[
 \left(\frac fg\right)^{(n)}=\frac{h_n}{g^{2^n}}.
\]
For $n=0$, choose $h_0=f$; the denominator is $g^{2^0}=g$,
so the formula is correct.
\step{Differentiate the inductive expression.}
Suppose the formula holds for some $n$, and write $m=2^n$.
The function $g^m$ is smooth, by the product rule applied a finite
number of times, and it never vanishes. Thus the ordinary quotient
rule, requiring only first derivatives at this step, gives
\[
 \left(\frac{h_n}{g^m}\right)'
 =\frac{h_n'g^m-h_n\,m g^{m-1}g'}{g^{2m}}.
\]
Define
\[
 h_{n+1}=h_n'g^m-mh_n g^{m-1}g'.
\]
All factors are smooth, so finite sums and products show that
$h_{n+1}$ is smooth. Since $2m=2^{n+1}$, this is the desired
formula at the next order. The induction also proves the existence
of that next derivative.
\step{Check continuity, not only existence.}
For each $n$, the formula expresses the $n$th derivative as a
quotient of continuous functions with a denominator that is nowhere
zero. That derivative is therefore continuous. We have proved the
existence and continuity of derivatives of every order.
\answer{$f/g$ is infinitely differentiable. The nonvanishing
assumption is what allows every quotient-rule step.}
